\documentclass[11pt,a4paper]{article}
\usepackage[margin=1in]{geometry}
\usepackage{amsmath,amssymb,amsthm}
\usepackage{algorithm}
\usepackage{algpseudocode}
\usepackage{booktabs}
\usepackage{hyperref}

\theoremstyle{definition}
\newtheorem{definition}{\textbf{Definition}}[section]
\newtheorem{theorem}{\textbf{Theorem}}[section]
\newtheorem{lemma}[theorem]{\textbf{Lemma}}
\newtheorem{remark}{\textbf{Remark}}[section]

\title{\textbf{Normalizing Flows}}
\author{\textbf{Team Scaibu} \\ \small Email: \texttt{jaydeep.9664920749@gmail.com} \\ \small \textbf{Architecture and Core Engine Team}}
\date{\textbf{October 2026}}

\begin{document}
\maketitle

\section{Definitions and Cognitive Mental Models}

\subsection{Formal Mathematical Definition}
\textbf{Definition (Normalizing Flow Operator):}
Let $\mathcal{Z} = \mathbb{R}^D$ denote a latent base space equipped with an analytically tractable probability density $p_z(\mathbf{z}) = \mathcal{N}(\mathbf{0}, \mathbf{I}_D)$ and let $\mathcal{X} = \mathbb{R}^D$ denote the observational data space. A \textbf{Normalizing Flow} is a diffeomorphism $f_\theta: \mathcal{Z} \to \mathcal{X}$ formed by the sequential composition of $K$ invertible, continuously differentiable transformations:
{\boldmath
\begin{equation}
\mathbf{x} = f_\theta(\mathbf{z}) = f_K \circ f_{K-1} \circ \dots \circ f_1(\mathbf{z}), \qquad \mathbf{z} = f_\theta^{-1}(\mathbf{x}) = f_1^{-1} \circ \dots \circ f_K^{-1}(\mathbf{x})
\end{equation}
}
By the multivariable \textbf{Change of Variables Theorem}, the exact marginal log-likelihood of observation $\mathbf{x}$ is evaluated without Monte Carlo approximation or lower-bound relaxation:
{\boldmath
\begin{equation}
\log p_\theta(\mathbf{x}) = \log p_z(f_\theta^{-1}(\mathbf{x})) - \log \left| \det \mathbf{J}_{f_\theta}(f_\theta^{-1}(\mathbf{x})) \right| = \log p_z(\mathbf{z}) + \sum_{k=1}^K \log \left| \det \mathbf{J}_{f_k^{-1}}(\mathbf{x}_{k}) \right|
\end{equation}
}
where $\mathbf{J}_f(\mathbf{z}) = \left[ \frac{\partial f_i}{\partial z_j} \right]_{i, j=1}^D$ is the Jacobian matrix. In \textbf{Affine Coupling Layers} (RealNVP), given partition boundary $d < D$:
{\boldmath
\begin{equation}
\mathbf{x}_{1:d} = \mathbf{z}_{1:d}, \qquad \mathbf{x}_{d+1:D} = \mathbf{z}_{d+1:D} \odot \exp(s(\mathbf{z}_{1:d})) + t(\mathbf{z}_{1:d})
\end{equation}
}
where scale $s$ and translation $t$ are arbitrary non-linear neural networks whose Jacobians never need to be computed or inverted.

\subsection{Expanded Non-Technical Definition (Intuition for Anyone)}
\textbf{Intuitive Conceptual Definition:}
Most generative models have a major flaw: GANs cannot compute exact probabilities, and VAEs only approximate probabilities using an imprecise lower bound (ELBO).

Normalizing Flows solve this by using purely **two-way invertible mathematical tunnels**:
\begin{enumerate}
    \item \textbf{The Invertible Pipe}: Imagine a flexible sheet of rubber. A flow starts with a simple standard bell curve (easy to sample from) and stretches, twists, and bends it into the complex shape of real data.
    \item \textbf{Exact Reversibility}: Because every stretching step is strictly reversible, you can take any real image and run it backward through the tunnel to find its exact original noise coordinates.
    \item \textbf{Volume Accounting (The Jacobian)}: Whenever you stretch or squeeze a sheet of rubber, its local thickness changes. The Jacobian determinant acts like a meticulous volume accountant that keeps track of exactly how much the space was stretched at every point.
    \item \textbf{Exact Density}: By combining the starting probability with the volume change, you compute the exact, 100\% true probability of any image.
\end{enumerate}

\subsection{Expanded Mental Model for Visualization (Understandable by a Child)}
\textbf{The Magic Origami Folder and Unfolder}:
Imagine a crisp, flat sheet of origami paper:
\begin{itemize}
    \item \textbf{The Smooth Paper Sheet (Base Distribution $\mathbf{z}$)}:
    You start with a simple square piece of paper. Anyone can draw a dot on it at random.
    
    \item \textbf{The Magic Folds (Flow Steps $f_k$)}:
    You fold the paper 10 times in very specific ways. Now the paper is a beautiful 3D origami swan ($\mathbf{x}$)!
    
    \item \textbf{The Reversible Unfolder ($f^{-1}$)}:
    Because you know exactly how each fold was made, you can unfold the swan step by step until it is a flat square sheet again—with zero tears or creases lost!
    
    \item \textbf{The Folding Magnifier (Jacobian Determinant)}:
    Some folds make the paper 4 layers thick; some folds unfold it flat. A tiny magnifying glass counts how many layers of paper are stacked under each point.
    
    \item \textbf{The Real vs Fake Detector}:
    If someone hands you a real swan folded with your rules, it unfolds into a clean square. If they hand you a crumpled ball of aluminum foil, the folding rules break down, giving you an impossible shape!
\end{itemize}

\section{Core Mathematical Equations: Formal Definitions, Meaning, and Importance}

\subsection{\textbf{Equation 1: Diffeomorphic Change of Variables Density}}
{\boldmath
\begin{equation}
p_X(\mathbf{x}) = p_Z(f^{-1}(\mathbf{x})) \cdot \left| \det \mathbf{J}_{f^{-1}}(\mathbf{x}) \right| = p_Z(\mathbf{z}) \cdot \left| \det \mathbf{J}_f(\mathbf{z}) \right|^{-1}
\end{equation}
}
\begin{itemize}
    \item \textbf{Formal Mathematical Definition}:
    The push-forward probability density transformation rule induced by smooth bijection $f$ with non-singular Jacobian.
    \item \textbf{What It Means}:
    The likelihood of data point $\mathbf{x}$ equals the prior probability of its pre-image $\mathbf{z}$ scaled by the local volume contraction factor.
    \item \textbf{Why It Is Important}:
    Allows exact likelihood computation and optimization without approximations, sampling variance, or intractable partition functions.
\end{itemize}

\subsection{\textbf{Equation 2: Affine Coupling Layer Forward Mapping}}
{\boldmath
\begin{equation}
\mathbf{x}_{1:d} = \mathbf{z}_{1:d}, \qquad \mathbf{x}_{d+1:D} = \mathbf{z}_{d+1:D} \odot \exp(s(\mathbf{z}_{1:d})) + t(\mathbf{z}_{1:d})
\end{equation}
}
\begin{itemize}
    \item \textbf{Formal Mathematical Definition}:
    The triangular diffeomorphism that acts as the identity on the first $d$ coordinates and applies coordinate-wise affine scaling and translation to the remaining $D - d$ coordinates conditioned on the first $d$.
    \item \textbf{What It Means}:
    Transforms half of the vector using complex neural networks $s$ and $t$ while keeping the other half as an immutable conditioning context.
    \item \textbf{Why It Is Important}:
    Guarantees that the Jacobian matrix is lower triangular, making its determinant trivial to compute in $\mathcal{O}(D)$ time rather than $\mathcal{O}(D^3)$.
\end{itemize}

\subsection{\textbf{Equation 3: Exact Analytical Inversion}}
{\boldmath
\begin{equation}
\mathbf{z}_{1:d} = \mathbf{x}_{1:d}, \qquad \mathbf{z}_{d+1:D} = (\mathbf{x}_{d+1:D} - t(\mathbf{x}_{1:d})) \odot \exp(-s(\mathbf{x}_{1:d}))
\end{equation}
}
\begin{itemize}
    \item \textbf{Formal Mathematical Definition}:
    The exact algebraic inverse of the affine coupling layer.
    \item \textbf{What It Means}:
    Inverts the transformation in a single feedforward pass using the exact same networks $s$ and $t$ without needing to invert $s$ or $t$ themselves.
    \item \textbf{Why It Is Important}:
    Allows $s$ and $t$ to be arbitrarily complex neural networks (e.g., ResNets, Convolutions, or MLPs) with zero invertibility constraints on their internal architectures.
\end{itemize}

\subsection{\textbf{Equation 4: Triangular Jacobian Log-Determinant}}
{\boldmath
\begin{equation}
\log \left| \det \mathbf{J}_f(\mathbf{z}) \right| = \log \left| \det \begin{bmatrix} \mathbf{I}_d & \mathbf{0} \\ \frac{\partial \mathbf{x}_{d+1:D}}{\partial \mathbf{z}_{1:d}} & \operatorname{diag}(\exp(s(\mathbf{z}_{1:d}))) \end{bmatrix} \right| = \sum_{i=1}^{D-d} s_i(\mathbf{z}_{1:d})
\end{equation}
}
\begin{itemize}
    \item \textbf{Formal Mathematical Definition}:
    The trace of the logarithm of the diagonal blocks of the triangular Jacobian matrix.
    \item \textbf{What It Means}:
    The total log volume expansion factor is simply the sum of the scale outputs of network $s$.
    \item \textbf{Why It Is Important}:
    Reduces an $\mathcal{O}(D^3)$ matrix determinant computation to a simple $\mathcal{O}(D)$ scalar sum, enabling real-time scaling to high-dimensional images.
\end{itemize}

\section{Rigorous Mathematical Analysis Checklist}

\subsection{[ ] Notation: Symbol and Operator Specification}
\begin{itemize}
    \item $\mathbf{z} \in \mathbb{R}^D$: Latent base Gaussian sample.
    \item $\mathbf{x} \in \mathbb{R}^D$: Ambient observation vector.
    \item $d \in \{1, \dots, D-1\}$: Split partition coordinate boundary.
    \item $s(\cdot) \in \mathbb{R}^{D-d}$: Scale transformation neural network.
    \item $t(\cdot) \in \mathbb{R}^{D-d}$: Translation transformation neural network.
    \item $\mathbf{J}_f$: $D \times D$ Jacobian matrix of partial derivatives.
    \item $\log p(\mathbf{x})$: Exact log-likelihood of observation $\mathbf{x}$.
\end{itemize}

\subsection{[ ] Definitions: Epistemic Nature of Variables}
\begin{table}[h]
\centering
\caption{\textbf{Epistemic Classification of Normalizing Flow Quantities}}
\begin{tabular}{lll}
\toprule
\textbf{Variable} & \textbf{Epistemic Classification} & \textbf{Mathematical Nature} \\
\midrule
$D, d$ & \textbf{Defined} (Dimension hyperparameters) & Natural numbers $\mathbb{N}_{\ge 1}$ \\
$\mathbf{W}_s, \mathbf{b}_s, \mathbf{W}_t, \mathbf{b}_t$ & \textbf{Learned} (Coupling network parameters) & Real weight matrices and vectors \\
$\mathbf{x}, \mathbf{z}$ & \textbf{Calculated} (Invertible representations) & Vectors in $\mathbb{R}^D$ \\
$\log |\det \mathbf{J}|, \log p(\mathbf{x})$ & \textbf{Calculated} (Exact density evaluations) & Real scalars in $\mathbb{R}$ \\
\bottomrule
\end{tabular}
\end{table}

\subsection{[ ] Operator Computation Graph}
{\boldmath
\begin{equation}
\mathbf{z} \xrightarrow{\text{Split}} (\mathbf{z}_1, \mathbf{z}_2) \xrightarrow{s(\mathbf{z}_1), t(\mathbf{z}_1)} \mathbf{x}_2 = \mathbf{z}_2 e^s + t \xrightarrow{\text{Concat}} \mathbf{x} \longrightarrow (\log p(\mathbf{z}), \sum s_i) \to \log p(\mathbf{x})
\end{equation}
}

\subsection{[ ] Dimensional Units and Tensor Ranks}
\begin{itemize}
    \item Input and output tensors $\mathbf{z}, \mathbf{x}$: Rank-1 tensors of shape $(D)$.
    \item Jacobian matrix $\mathbf{J}$: Rank-2 tensor of shape $(D, D)$.
    \item Log-determinant and log-likelihood: Dimensionless real scalars.
\end{itemize}

\subsection{[ ] Limiting Regimes and Asymptotic Behaviors}
\begin{itemize}
    \item \textbf{Zero Coupling Regime ($s \equiv 0, t \equiv 0$)}: Transformation is exact identity mapping $f(\mathbf{z}) = \mathbf{z}$; $\log |\det \mathbf{J}| = 0$.
    \item \textbf{Infinite Depth Limit ($K \to \infty$)}: Discrete flow composition converges to a continuous-time Neural Ordinary Differential Equation (Neural ODE).
\end{itemize}

\subsection{[ ] Operational Constraints and Failure Modes}
\begin{itemize}
    \item \textbf{Unbounded Scale Instability}: If $s(\mathbf{z}_1) \gg 0$, $\exp(s)$ causes floating-point overflow. Clamping or passing $s$ through bounded activations (e.g., $2 \tanh(s / 2)$) is mathematically mandatory.
\end{itemize}

\subsection{[ ] Mathematical Invariants and Conserved Quantities}
\begin{itemize}
    \item \textbf{Probability Mass Conservation}: $\int_{\mathcal{X}} p_X(\mathbf{x}) d\mathbf{x} = \int_{\mathcal{Z}} p_Z(\mathbf{z}) d\mathbf{z} \equiv 1.0$ holds analytically.
    \item \textbf{Exact Reversibility Invariant}: $\|f^{-1}(f(\mathbf{z})) - \mathbf{z}\| \le \epsilon_{\text{machine}}$.
\end{itemize}

\subsection{[ ] Cross-Domain Mathematical Isomorphisms}
\begin{itemize}
    \item \textbf{Symplectic Integrators in Hamiltonian Mechanics}: Volume-preserving flows preserve phase space volume, strictly mirroring Liouville's theorem in statistical mechanics.
\end{itemize}

\section{Theoretical Theorems and Formal Proofs}

\begin{theorem}[\textbf{Change of Variables for Probability Density Functions}]
Let $\mathbf{Z} \in \mathbb{R}^D$ be a random vector with continuous probability density $p_Z(\mathbf{z})$, and let $\mathbf{X} = f(\mathbf{Z})$ where $f: \mathbb{R}^D \to \mathbb{R}^D$ is a diffeomorphism (smooth bijection with smooth inverse). Then the probability density of $\mathbf{X}$ is given by:
{\boldmath
\begin{equation}
p_X(\mathbf{x}) = p_Z(f^{-1}(\mathbf{x})) \cdot \left| \det \left( \frac{\partial f^{-1}(\mathbf{x})}{\partial \mathbf{x}} \right) \right|
\end{equation}
}
\end{theorem}

\begin{proof}
Let $\mathcal{B} \subset \mathbb{R}^D$ be any arbitrary Borel-measurable set in data space $\mathcal{X}$, and let $f^{-1}(\mathcal{B}) \subset \mathcal{Z}$ denote its pre-image under $f$.
By probability conservation:
{\boldmath
\begin{equation}
\mathbb{P}(\mathbf{X} \in \mathcal{B}) = \mathbb{P}(\mathbf{Z} \in f^{-1}(\mathcal{B})) = \int_{f^{-1}(\mathcal{B})} p_Z(\mathbf{z}) \, d\mathbf{z}
\end{equation}
}
Apply the multivariable substitution rule of integration $\mathbf{z} = f^{-1}(\mathbf{x})$. The differential volume element transforms according to the Jacobian determinant:
{\boldmath
\begin{equation}
d\mathbf{z} = \left| \det \mathbf{J}_{f^{-1}}(\mathbf{x}) \right| d\mathbf{x}
\end{equation}
}
Substituting this into the integral:
{\boldmath
\begin{equation}
\int_{f^{-1}(\mathcal{B})} p_Z(\mathbf{z}) \, d\mathbf{z} = \int_{\mathcal{B}} p_Z(f^{-1}(\mathbf{x})) \left| \det \mathbf{J}_{f^{-1}}(\mathbf{x}) \right| d\mathbf{x}
\end{equation}
}
By definition of the probability density function $p_X(\mathbf{x})$:
{\boldmath
\begin{equation}
\mathbb{P}(\mathbf{X} \in \mathcal{B}) = \int_{\mathcal{B}} p_X(\mathbf{x}) \, d\mathbf{x}
\end{equation}
}
Since this equality holds for all measurable sets $\mathcal{B}$, the integrands must be equal almost everywhere:
{\boldmath
\begin{equation}
p_X(\mathbf{x}) = p_Z(f^{-1}(\mathbf{x})) \cdot \left| \det \mathbf{J}_{f^{-1}}(\mathbf{x}) \right|
\end{equation}
}
This completes the proof.
\end{proof}

\section{Foundational Architectural and Epistemic Meta-Questions}

\subsection{Architectural Question 1: How Do Coupling Layers Achieve Non-Linear Invertibility Without Inverting Their Neural Networks?}
\textbf{Question}: If arbitrary neural networks $s(\cdot)$ and $t(\cdot)$ are non-invertible, how does the overall coupling layer remain strictly invertible?

\textbf{Meta-Question}: Does partitioning a coordinate space allow decoupling functional evaluation from algebraic inversion?

\textbf{Answer}: In an affine coupling layer, the inputs to $s$ and $t$ are $\mathbf{z}_{1:d}$, which are copied unchanged to the output $\mathbf{x}_{1:d} = \mathbf{z}_{1:d}$. During inversion, the exact input $\mathbf{z}_{1:d}$ is already known because $\mathbf{x}_{1:d}$ is directly available. Therefore, the inverse step simply evaluates $s(\mathbf{x}_{1:d})$ and $t(\mathbf{x}_{1:d})$ in the forward direction. Inversion only requires subtracting $t$ and dividing by $\exp(s)$, which are simple scalar operations. Networks $s$ and $t$ never need to be run backward or inverted.

\section{Algorithmic Specification}

\begin{algorithm}[H]
\caption{Affine Coupling Normalizing Flow Forward and Inverse Evaluation}
\label{alg:norm_flow}
\begin{algorithmic}[1]
\Require Vector $\mathbf{u} \in \mathbb{R}^D$, Weights $\mathbf{W}_s, \mathbf{b}_s, \mathbf{W}_t, \mathbf{b}_t$, Split $d$, Mode $\text{inverse} \in \{\text{true}, \text{false}\}$
\Ensure Output vector $\mathbf{v} \in \mathbb{R}^D$, Log-det Jacobian $\mathcal{J}$, Log-likelihood $\mathcal{L}$
\State $\mathbf{u}_1 \gets \mathbf{u}[1 \dots d], \quad \mathbf{u}_2 \gets \mathbf{u}[d+1 \dots D]$
\State $D_{\text{rem}} \gets D - d$
\State $\mathbf{s} \gets \text{zeros}(D_{\text{rem}}), \quad \mathbf{t} \gets \text{zeros}(D_{\text{rem}})$
\For{$i \gets 1 \textbf{ to } D_{\text{rem}}$}
    \State $s_{\text{act}} \gets b_{s, i} + \sum_{j=1}^d W_{s, i, j} \cdot u_{1, j}$
    \State $t_{\text{act}} \gets b_{t, i} + \sum_{j=1}^d W_{t, i, j} \cdot u_{1, j}$
    \State $s_i \gets 2.0 \cdot \tanh(0.5 \cdot s_{\text{act}})$ \Comment{Clamped scale}
    \State $t_i \gets t_{\text{act}}$
\EndFor
\State $\mathbf{v} \gets \text{copy}(\mathbf{u}_1)$
\If{$\text{inverse} = \text{false}$}
    \For{$i \gets 1 \textbf{ to } D_{\text{rem}}$}
        \State Append $u_{2, i} \cdot \exp(s_i) + t_i$ to $\mathbf{v}$
    \EndFor
    \State $\mathcal{J} \gets \sum_{i=1}^{D_{\text{rem}}} s_i$
\Else
    \For{$i \gets 1 \textbf{ to } D_{\text{rem}}$}
        \State Append $(u_{2, i} - t_i) \cdot \exp(-s_i)$ to $\mathbf{v}$
    \EndFor
    \State $\mathcal{J} \gets -\sum_{i=1}^{D_{\text{rem}}} s_i$
\EndIf
\State $\mathbf{z}_{\text{eval}} \gets \mathbf{u} \text{ if not inverse else } \mathbf{v}$
\State $\text{sum\_sq} \gets \sum_{k=1}^D z_{\text{eval}, k}^2$
\State $\log p_z \gets -0.5 \cdot (D \cdot \ln(2\pi) + \text{sum\_sq})$
\State $\mathcal{L} \gets \log p_z - \mathcal{J}$
\State \Return $\{\text{output}: \mathbf{v}, \text{log\_det\_jacobian}: \mathcal{J}, \text{log\_likelihood}: \mathcal{L}\}$
\end{algorithmic}
\end{algorithm}

\section{Complexity Analysis}
\begin{itemize}
    \item \textbf{Time Complexity}:
    \begin{itemize}
        \item Scale and translation evaluation: $2(D - d) \times d$ operations ($\Theta((D - d)d)$).
        \item Coordinate affine transformation: $D - d$ operations ($\Theta(D - d)$).
        \item Total Time: $\Theta((D - d)d)$ FLOPs.
    \end{itemize}
    \item \textbf{Space Complexity}:
    \begin{itemize}
        \item Output vector $\mathbf{v}$: $\Theta(D)$ floats.
        \item Auxiliary scale and translation buffers: $\Theta(D - d)$ floats.
        \item Total Space: $\Theta(D)$.
    \end{itemize}
\end{itemize}

\section{Repository Reference}
All mathematical formulations, algorithmic implementations, contracts, and test suites are maintained at:
\begin{center}
\url{https://github.com/Chief-Strategist-J/llm-obs-infra}
\end{center}

\end{document}
